% ID: SEL-D0063-OA-E02
% ESTADO: validada. S063.
\textbf{a)} La condición de extremo interior exige $f'(-2)=0$. Como $f'(x)=4x+a$, queda $a=8$. Además, $f(1)=2+a+b=3$, luego $\boxed{b=-7}$. El extremo es un mínimo porque la segunda derivada es $4>0$.

\textbf{b)} Con $a=8$ y $b=-10$, la parábola es convexa en todo $\mathbb R$, pues $f''(x)=4$. Al completar el cuadrado se obtiene $f(x)=2(x+2)^2-18$. El mínimo es $\boxed{-18}$ en $x=-2$. La función se anula cuando $(x+2)^2=9$, es decir, en $\boxed{x=-5,\ 1}$.
